%=============================================================================!
% 1.9  FROM BALLS TO ATOMS                                                    !
%=============================================================================!
\section{From balls to atoms}
\label{sec:mo-atoms}

None of the preceding sections depended on the moving object being a
ball, and the same description applies directly to atoms. A system of $N$
atoms has a position vector for each atom, $\vb{r}_1(t), \dots,
\vb{r}_N(t)$, and each atom has its own velocity $\vb{v}_i =
\dd\vb{r}_i/\dd t$ and acceleration $\vb{a}_i = \dd\vb{v}_i/\dd t$. When
all the positions are collected into a single list of $3N$ numbers,
$\vb{r}^N = (\vb{r}_1, \dots, \vb{r}_N)$ (the superscript is a label,
not a power), the motion of the whole system becomes a single path
$\vb{r}^N(t)$: just as three numbers locate a point in space, $3N$
numbers locate a point in a space of $3N$ dimensions, which cannot be
drawn but obeys the same component-by-component rules, and every result of the preceding sections applies to each of its components.
The only change is in the units, which become \si{\angstrom},
\si{\femto\second}, \si{\angstrom\per\femto\second} and
\si{\angstrom\per\femto\second\squared}.

\subsection*{Velocities from frames}

A molecular dynamics calculation does not record this path continuously.
It stores frames instead, the positions of all the atoms at times
separated by an interval $\tau$, and the velocities must sometimes be
recovered from these frames afterwards. The derivative, a limit as $\tau
\to 0$, must then be replaced by a finite difference with $\tau$ fixed, and
the Taylor series of \cref{eq:mo-taylor} tells us how large an error this
introduces. We treat one coordinate $x$; every component of every atom
behaves in the same way.

The simplest choice is the forward difference, the secant slope of
\cref{sec:mo-velocity}. Subtracting $x(t)$ from both sides of
\cref{eq:mo-taylor} and dividing by $\tau$,
\begin{equation*}
   \frac{x(t+\tau) - x(t)}{\tau} = v(t) + \tfrac12 a(t)\,\tau
   + \order{\tau^2},
\end{equation*}
where the term $\tfrac16\dddot{x}\,\tau^2$ and the remainder, now
$\order{\tau^4}/\tau = \order{\tau^3}$, are both no larger than a constant
times $\tau^2$ when $\tau$ is small, and so together are $\order{\tau^2}$.
So the forward difference returns the velocity with an error whose leading
term, $\tfrac12 a\tau$, is proportional to $\tau$. For the ball this is
the term $-\tfrac12 g\tau$ of \cref{eq:mo-secant}.

A better estimate uses the frames on both sides of $t$. Writing the
Taylor series forwards and backwards in time, by replacing $\tau$ with
$-\tau$ in the second,
\begin{align*}
   x(t+\tau) &= x + v\tau + \tfrac12 a\tau^2 + \tfrac16\dddot{x}\tau^3
      + \tfrac1{24}x^{(4)}\tau^4 + \order{\tau^5},\\
   x(t-\tau) &= x - v\tau + \tfrac12 a\tau^2 - \tfrac16\dddot{x}\tau^3
      + \tfrac1{24}x^{(4)}\tau^4 + \order{\tau^5},
\end{align*}
where $x$, $v$, $a$, $\dddot{x}$ and $x^{(4)}$ are evaluated at $t$. The terms with
even powers of $\tau$ have the same sign in both lines, and those with odd
powers have opposite signs. Subtracting the second line from the first
therefore cancels the even terms and doubles the odd ones,
\begin{equation*}
   x(t+\tau) - x(t-\tau) = 2v\tau + \tfrac13\dddot{x}\,\tau^3
   + \order{\tau^5},
\end{equation*}
where the remainder is of order $\tau^5$ because the $\tau^4$ terms also
cancel. Dividing by $2\tau$ gives the central difference,
\begin{equation*}
   \frac{x(t+\tau) - x(t-\tau)}{2\tau} = v(t) + \tfrac16\,\dddot{x}(t)\,
   \tau^2 + \order{\tau^4},
\end{equation*}
whose error is proportional to $\tau^2$. Halving the interval therefore
halves the error of the forward difference but quarters that of the
central difference. In \texttt{mdlab} the central difference occupies a
single line.

\begin{code}
def central_difference(positions: ArrayLike, tau: float) -> NDArray:
    x = np.asarray(positions, dtype=float)
    return (x[2:] - x[:-2]) / (2.0 * tau)
\end{code}

Listings in this book are taken from \texttt{mdlab} without their
docstrings, which give the shape and unit of every argument. Python
counts from 0, so \texttt{x[2:]} holds the frames from the third onwards
and \texttt{x[:-2]} those up to the third from last; entry $k$ of their
difference is $x_{k+2} - x_k$, the central difference at frame $k + 1$,
so the first and last frames receive no velocity. \texttt{np.asarray}
turns the input into a NumPy array of decimal numbers, and the
annotations after the colons and the arrow only document the types
expected and returned.

\subsection*{A vibrating coordinate}

The atoms in molecules and solids vibrate (Chapter~4), and a vibrating
coordinate can be modelled by the sinusoid of \cref{sec:mo-trig} with zero
phase, $x(t) = A\sin\omega t$, of amplitude $A$, angular frequency $\omega$
and period $2\pi/\omega$. For this motion the central difference
can be evaluated exactly. Subtracting the two sine formulae of
\cref{eq:mo-addition},
\begin{align*}
   \sin(p+q) - \sin(p-q)
   &= (\sin p\cos q + \cos p\sin q) - (\sin p\cos q - \cos p\sin q)\\
   &= 2\cos p\sin q .
\end{align*}
With $p = \omega t$ and $q = \omega\tau$, and multiplying through by
$A/2\tau$,
\begin{equation}
   \frac{x(t+\tau) - x(t-\tau)}{2\tau}
   = \frac{A\cos\omega t\,\sin\omega\tau}{\tau}
   = \underbrace{A\omega\cos\omega t}_{v(t)}\;
     \frac{\sin\omega\tau}{\omega\tau} ,
   \label{eq:mo-sinc}
\end{equation}
where the last step multiplies and divides by $\omega$ so that the true
velocity $v(t) = A\omega\cos\omega t$ appears. Every recovered velocity is
therefore too small by the same factor, $\sin(\omega\tau)/(\omega\tau)$.
To see how close this factor is to 1, we divide the series for $\sin y$ in
\cref{eq:mo-sin-series} by $y$,
\begin{equation*}
   \frac{\sin y}{y} = 1 - \frac{y^2}{6} + \order{y^4} .
\end{equation*}
The relative error, the error divided by the true velocity, is therefore
\begin{equation*}
   \frac{v(t) - v(t)\sin(\omega\tau)/(\omega\tau)}{v(t)}
   = 1 - \frac{\sin\omega\tau}{\omega\tau} \approx \frac{(\omega\tau)^2}{6} .
\end{equation*}

For an amplitude of \SI{0.1}{\angstrom} and a period of
\SI{20}{\femto\second}, the angular frequency is $\omega =
2\pi/\SI{20}{\femto\second} = \SI{0.3142}{\per\femto\second}$. Frames
saved every \SI{1}{\femto\second} then give $\omega\tau = 0.3142$ and an
error of $0.3142^2/6 = \SI{1.645}{\percent}$ by the approximation, against
\SI{1.637}{\percent} exactly. Frames every \SI{0.1}{\femto\second} give
\SI{0.016}{\percent} either way, and frames every \SI{2}{\femto\second}
give \SI{6.58}{\percent} by the approximation against \SI{6.45}{\percent}
exactly (\cref{fig:mo-differences}). When $\tau$ equals half a
period, $\omega\tau = \pi$ and the factor $\sin\pi/\pi$ vanishes
altogether: $x(t+\tau)$ and $x(t-\tau)$ are then a whole period apart and
therefore equal, and the central difference detects no motion at all.

\begin{figure}[tb]
\centering
\includegraphics{ch01_motion/differences}
\caption{Velocities recovered from frames. A coordinate $x = A\sin\omega
t$ with $A = \SI{0.1}{\angstrom}$ and period $2\pi/\omega =
\SI{20}{\femto\second}$ is sampled every $\tau$. (a) Largest error of the
forward and central differences, with their leading terms
$A\omega^2\tau/2$ and $A\omega^3\tau^2/6$ (dashed). (b) Relative error of
the central difference, $1 - \sin(\omega\tau)/(\omega\tau)$, against its
leading term; the dotted line marks $\tau = \SI{1}{\femto\second}$.}
\label{fig:mo-differences}
\end{figure}

\begin{practice}
Frames must be closely spaced compared with the fastest motion that they
are intended to resolve, and \cref{eq:mo-sinc} makes this requirement
quantitative. The coordinates of real atoms move as sums of sinusoids of
several frequencies (Chapter~4). The central difference of a sum is the
sum of the central differences of its terms, so it shrinks the velocity
of each sinusoid by that sinusoid's own factor $\sin(\omega\tau)/
(\omega\tau)$. If the interval satisfies $(\omega_{\max}\tau)^2/6 \le
\epsilon$ for the highest angular frequency $\omega_{\max}$ present, each
sinusoid's velocity is recovered to a relative accuracy $\epsilon$, and
the error in the whole velocity is at most $\epsilon$ times the sum of
the largest speeds $A\omega$ of the sinusoids, since the size of a sum is
at most the sum of the sizes of its terms. This bounds the error
relative to the typical speed of the motion, not relative to the velocity
at each instant: where the sinusoids cancel and the velocity passes
through zero, even a small error is a large fraction of it. For $x =
\sin t - \tfrac12\sin 2t$, in metres with $t$ in seconds, and $\tau =
\SI{0.1}{\second}$, $(\omega_{\max}\tau)^2/6 = \SI{0.67}{\percent}$ with
$\omega_{\max} = 2$. The true velocity is $\cos t - \cos 2t$, while by
\cref{eq:mo-sinc} applied to each term the central difference is $\cos
t\,(\sin 0.1/0.1) - \cos 2t\,(\sin 0.2/0.2)$. At $t =
\SI{0.01}{\second}$, where the true velocity is
\SI{1.5e-4}{\metre\per\second}, the central difference gives
\SI{5.1e-3}{\metre\per\second}, 34 times too large. A path saved less
frequently than the criterion allows should store the velocities
directly. A related requirement limits the
size of the step that molecular dynamics itself can take (Chapter~9).
\end{practice}

\notebook{01}{sample an oscillation and recover its velocity}
